
Pour le premier il s'agit de transformer @@M0@@ en la matrice unité @@M1@@ dans un voisinage zariskien de @@M2@@ en multipliant @@M3@@ à gauche et à droite par des matrices convenables @@M4@@. Montrons d'abord qu'on peut rendre ainsi @@M5@@. Si @@M6@@ et @@M7@@, nous cherchons une matrice @@M8@@ telle que le produit @@M9@@ ait @@M10@@. Puisque @@M11@@, il existe au moins un indice @@M12@@ tel que @@M13@@. Alors on a @@M14@@, et @@M15@@ répond à la question. On peut donc supposer @@M16@@, et quitte à localiser, on peut supposer @@M17@@ inversible. Si @@M18@@, alors dans @@M19@@ on a
@@M20@@. Quitte à localiser, on peut alors supposer @@M21@@ inversible. Alors @@M22@@ est inversible et on trouve @@M23@@ dans @@M24@@.

Supposant @@M25@@ on peut alors rendre @@M26@@ à la forme où
@@M27@@
dans un voisinage zariskien, de @@M28@@ en multipliant @@M29@@ à gauche (resp. à droite) par des matrices @@M30@@ (resp. @@M31@@). Puis on va rendre @@M32@@ etc. Par récurrence on peut alors transformer @@M33@@ dans un voisinage zariskien @@M34@@ de @@M35@@ en une matrice @@M36@@ ayant @@M37@@ pour @@M38@@ ou @@M39@@. Mais par construction même on a alors @@M40@@ sur @@M41@@, où @@M42@@ et @@M43@@ sont des produits de matrices @@M44@@. Donc on a
@@M45@@
de sorte que @@M46@@, @@M47@@, et on a gagné.

On suppose donc que @@M48@@ soit déjà sur @@M49@@ un produit de matrices @@M50@@.
Reste à voir que celles-ci sont des produits de commutateurs dans un voisinage étale de @@M51@@. Tout d'abord on voit que deux matrices
